\documentclass[../../main.tex]{subfiles}
\begin{document}
\section*{习题}\label{section:ch12-习题}

\begin{exercise}\label{exercise:ch12-01}
给出有限维向量空间 $V$ 与 $W$ 以及一个具体元素 $\alpha\in V\otimes W$，使它不能写成 $v\otimes w$（$v\in V$，$w\in W$）。
\end{exercise}

\begin{exercise}\label{exercise:ch12-02}
对任意有限维实向量空间 $V$，证明存在典范同构 $\mathbb{R}\otimes V\cong V\cong V\otimes\mathbb{R}$。
\end{exercise}

\begin{exercise}\label{exercise:ch12-03}
设 $V$ 与 $W$ 为有限维实向量空间。证明张量积空间 $V\otimes W$ 在典范同构意义下由其特征性质唯一决定。更精确地说，设 $\widetilde\pi\colon V\times W\to Z$ 是到向量空间 $Z$ 的双线性映射，满足：对任意双线性映射 $B\colon V\times W\to Y$，存在唯一线性映射 $\widetilde B\colon Z\to Y$ 使下图交换。则存在唯一同构 $\Phi\colon V\otimes W\to Z$ 使 $\widetilde\pi=\Phi\circ\pi$，其中 $\pi\colon V\times W\to V\otimes W$ 是典范投影。（注：这表明用来定义张量积空间的具体构造细节无关紧要，只要所得空间满足特征性质即可。）
\end{exercise}

\begin{exercise}\label{exercise:ch12-04}
设 $V_1,\dots,V_k$ 与 $W$ 为有限维实向量空间。证明存在典范（与基无关）同构
\begin{align*}
V_1^*\otimes\cdots\otimes V_k^*\otimes W\cong L(V_1,\dots,V_k;W).
\end{align*}
（特别地，这意味着 $V^*\otimes W$ 典范同构于从 $V$ 到 $W$ 的线性映射空间 $L(V;W)$。）
\end{exercise}

\begin{exercise}\label{exercise:ch12-05}
设 $V$ 为 $n$ 维实向量空间。证明
\begin{align*}
\dim\Sigma^k(V^*)=\binom{n+k-1}{k}=\frac{(n+k-1)!}{k!(n-1)!}.
\end{align*}
\end{exercise}

\begin{exercise}\label{exercise:ch12-06}
(a) 设 $\alpha$ 是有限维实向量空间 $V$ 上的协变 $k$ 张量。证明 $\mathrm{Sym}\,\alpha$ 是唯一满足 $(\mathrm{Sym}\,\alpha)(v,\dots,v)=\alpha(v,\dots,v)$（对所有 $v\in V$）的对称 $k$ 张量。

(b) 证明对称积结合：对所有对称张量 $\alpha,\beta,\gamma$，有 $(\alpha\beta)\gamma=\alpha(\beta\gamma)$。（提示：用 (a)。）

(c) 设 $\omega^1,\dots,\omega^k$ 是有限维向量空间上的余向量。证明它们的对称积满足
\begin{align*}
\omega^1\cdots\omega^k=\frac{1}{k!}\sum_{\sigma\in S_k}\omega^{\sigma(1)}\otimes\cdots\otimes\omega^{\sigma(k)}.
\end{align*}
\end{exercise}

\begin{exercise}\label{exercise:ch12-07}
设 $(\varepsilon^1,\varepsilon^2,\varepsilon^3)$ 是 $(\mathbb{R}^3)^*$ 的标准对偶基。证明 $\varepsilon^1\otimes\varepsilon^2\otimes\varepsilon^3$ 不能写成一个交错张量与一个对称张量之和。
\end{exercise}

\begin{exercise}\label{exercise:ch12-08}
设 $M$ 为光滑 $n$ 流形，$A$ 是 $M$ 上的光滑协变 $k$ 张量场。若 $(U,(x^i))$ 与 $(\widetilde U,(\widetilde x^j))$ 是 $M$ 上两个相交的光滑坐标图，则可写
\begin{align*}
A=A_{i_1\dots i_k}\,dx^{i_1}\otimes\cdots\otimes dx^{i_k}
=\widetilde A_{j_1\dots j_k}\,d\widetilde x^{j_1}\otimes\cdots\otimes d\widetilde x^{j_k}.
\end{align*}
求一个类似于 (11.7) 的变换律，用 $\widetilde A_{j_1\dots j_k}$ 表出分量 $A_{i_1\dots i_k}$。
\end{exercise}

\begin{exercise}\label{exercise:ch12-09}
把习题~12-8~的坐标变换律推广到任意型的混合张量场。
\end{exercise}

\begin{exercise}\label{exercise:ch12-10}
证明对每对非负整数 $k,l$ 与每个微分同胚 $F\colon M\to N$，存在前推与拉回同构
\begin{align*}
F_*\colon\Gamma(T^{(k,l)}TM)\to\Gamma(T^{(k,l)}TN),\qquad
F^*\colon\Gamma(T^{(k,l)}TN)\to\Gamma(T^{(k,l)}TM),
\end{align*}
使得 $F^*$ 在协变张量场上与通常拉回一致，$F_*$ 在反变 $1$ 张量场（即向量场）上与通常前推一致，且满足：
\begin{enumerate}[(a)]
\item $F_*=(F^*)^{-1}$；
\item $F^*(A\otimes B)=F^*A\otimes F^*B$；
\item $(F\circ G)_*=F_*\circ G_*$；
\item $(F\circ G)^*=G^*\circ F^*$；
\item $(\mathrm{Id}_M)_*=(\mathrm{Id}_M)^*=\mathrm{Id}$；
\item $F^*(A(X_1,\dots,X_k))=F^*A(F_*^{-1}(X_1),\dots,F_*^{-1}(X_k))$，其中 $A\in\mathcal T^k(N)$，$X_1,\dots,X_k\in\mathfrak X(N)$。
\end{enumerate}
\end{exercise}

\begin{exercise}\label{exercise:ch12-11}
设 $M$ 为光滑流形，$A$ 是 $M$ 上的光滑协变张量场，$V,W\in\mathfrak X(M)$。证明
\begin{align*}
\mathcal L_V\mathcal L_W A-\mathcal L_W\mathcal L_V A=\mathcal L_{[V,W]}A.
\end{align*}
（提示：对 $A$ 的秩作归纳，从 $1$ 张量场的推论~12.33~开始。）
\end{exercise}

\begin{exercise}\label{exercise:ch12-12}
设 $M$ 为光滑流形，$V\in\mathfrak X(M)$。证明协变张量场上的 Lie 导数算子 $\mathcal L_V\colon\mathcal T^k(M)\to\mathcal T^k(M)$（$k\geqslant 0$）由下列性质唯一刻画：
\begin{enumerate}[(a)]
\item $\mathcal L_V$ 对 $\mathbb{R}$ 线性；
\item $\mathcal L_V f=Vf$，对 $f\in\mathcal T^0(M)=C^\infty(M)$；
\item $\mathcal L_V(A\otimes B)=\mathcal L_V A\otimes B+A\otimes\mathcal L_V B$，对 $A\in\mathcal T^k(M)$，$B\in\mathcal T^l(M)$；
\item $\mathcal L_V(\omega(X))=(\mathcal L_V\omega)(X)+\omega([V,X])$，对 $\omega\in\mathcal T^1(M)$，$X\in\mathfrak X(M)$。
\end{enumerate}
（注：张量场上的 Lie 导数算子有时被定义为满足这些性质的唯一算子。这样定义的优点是在带边流形上也有意义——在那里 $V$ 的流可能不存在。）
\end{exercise}
\end{document}
